\documentclass{ximera}
%verbeterd 13/11/26
\title{Stand van zaken van de cursus}
\author{Daan Lipkens}

\begin{document}
\begin{abstract}
    Een overzicht van de afgewerkte en nagekeken delen.
\end{abstract}
\maketitle

De cursus wordt in eerste instantie geschreven voor de wetenschapsrichtingen. Daarna pas ik de inhoud aan voor jullie, die maar één uur fysica hebben.

\section*{Geschreven}
\subsection*{Geschreven voor 1 uur fysica}
\begin{itemize}
    \item Hoofdstuk 1: elektrische krachtwerking
    \item Hoofdstuk 2: elektrisch veld
    \item Hoofdstuk 4: gemengde schakelingen
    \item Hoofdstuk 5: het magnetisch veld
    \item Appendix: vectorrekenen
    \item Appendix: arbeid en energie
    \item Appendix: natuurconstanten en gegevens
    \item Appendix: standaardvorm
\end{itemize}

\subsection*{Geschreven voor 2 uur fysica}
\begin{itemize}
    \item Hoofdstuk 1: elektrische krachtwerking
    \item Hoofdstuk 2: elektrisch veld
    \item Hoofdstuk 3: potentiaal
    \item Hoofdstuk 4: het magnetisch veld
    \item Appendix: standaardvorm
    \item Appendix: vectorrekenen
    \item Appendix: arbeid en energie
    \item Appendix: natuurconstanten en gegevens
\end{itemize}

\section*{Nagekeken}
\begin{itemize}
    \item Hoofdstuk 1: elektrische krachtwerking
    \item Hoofdstuk 2: elektrisch veld
\end{itemize}

\section*{Wat komt er nu?}
\begin{itemize}
    \item Appendix: eenheden en standaardvorm (herschreven) nalezen
    \item Hoofdstuk 3: potentiaal nalezen en aanpassen voor 1 uur fysica
    \item Hoofdstuk 3: oefeningen voor 1 uur fysica nakijken
    \item Hoofdstuk 4: gemengde schakelingen nakijken
    \item Hoofdstuk 5: het magnetisch veld nalezen
    \item Herhalingsoefeningen per trimester
    \item Oefeningen van toelatingsproeven
\end{itemize}

\section*{Feedback}
Is er iets niet duidelijk, staat er ergens een fout of mis je iets? Laat het dan zeker weten.

\end{document}